\section{Related Work}
\label{sec:related}
\subsection{Audiovisual Conversation and Response Generation}
% Paragraph role P01: 视觉倾听沿革
Listening-generation methods synthesize facial expressions and head movements from a speaker's communication signals. ViCo~\cite{vico} introduces responsive listening-head generation and paired data, using speaker audiovisual signals and a preset attitude to generate listener motion. Learning to Listen~\cite{learningtolisten} models listener motion as a nondeterministic distribution, allowing multiple responses to the same input. Subsequent work explores diffusion-based diverse generation, emotion control, and text-based customization through MFR-Net~\cite{mfrnet}, ELP~\cite{elp}, and CustomListener~\cite{customlistener}, respectively. DiffListener~\cite{difflistener} and DiTaiListener~\cite{ditailistener} further develop expressive listening-video generation.

% Paragraph role P02: 行为适当性与事件表示
Beyond visual quality, ReactFace~\cite{reactface} studies online generation of multiple appropriate facial reactions, considering appropriateness, diversity, and synchrony. GLARE~\cite{glare} represents reactions as events with categories and temporal intervals, combining prosodic conditioning, temporal reaction supervision, and event evaluation. These visual listening methods provide a basis for behavior modeling, but do not generate the listener's own audio. Our task includes brief acknowledgments and laughter together with the corresponding video, while retaining silent listening as a valid output.

% Paragraph role P03: 已知音轨的交互视频
Interactive-avatar methods also differ in whether the target audio is supplied or generated. INFP~\cite{infp} and DyStream~\cite{dystream} use both participants' audio tracks to generate avatar videos that include speaking and listening. ARIG~\cite{arig} additionally models interaction using both participants' audiovisual signals and dialogue states. These methods incorporate role changes and listening into generation, but receive the target person's audio as a condition. Our model must generate both listener audio and video from the other person's current speech and a listener reference frame.

% Paragraph role P04: 自主AV回应与本文边界
OmniResponse~\cite{omniresponse} is closer to this setting: it uses speaker facial information, transcripts, and dialogue history to generate aligned text and facial responses, with TempoVoice synthesizing audio for vocal responses. OmniPersona~\cite{omnipersona} encodes multimodal conversation history into linguistic and non-linguistic conditions and jointly generates audio and video, including listening backchannels and identity consistency across turns. Audiovisual listening responses therefore already have a research basis. Our focus is feedback during the other person's turn from a short current audio segment and a listener first frame, without dialogue history, speaker video, a target listener audio track, or a reference voice. We address this setting through direct audio conditioning, target feature supervision, and clip-level event supervision, rather than introducing a multi-turn memory task.

\subsection{MLLM-Based Cascaded Generation}
\cnnote{中文待补：作者已确认的新结构}{本小节尚无写定正文，先保留中文。后续结合OmniPersona对理解与生成框架的介绍，核实“MLLM理解／规划→文本或隐表征→生成器”的代表工作、接口与条件，再写文献介绍。不能把所有条件接口统称为同一种级联；Qwen→OVI系统对照也不等于先生成音频再驱动视频。这里不新增“联合生成优于级联”的主张，暂不补造文献或技术关系。}

\subsection{Conversation Data and Feedback Evaluation}
% Paragraph role P05: 数据资源及本文组织
Existing datasets cover different aspects of dyadic interaction. ViCo~\cite{vico} provides listening-generation samples. RealTalk~\cite{realtalk} contains natural dyadic conversations with transcripts and active-speaker annotations, while Seamless Interaction~\cite{seamless} provides large-scale face-to-face interaction data. SpeakerVid-5M~\cite{speakervid} organizes dialogue, listening, and multi-turn branches; its listening branch pairs speaker audio with silent listener video. ResponseNet, introduced with OmniResponse~\cite{omniresponse}, provides synchronized videos of both participants, multichannel audio, transcripts, and facial behavior annotations. GLARE~\cite{glare} adds visual reaction categories and temporal intervals. These resources already provide paired media and several types of behavior annotation. Their support for our task depends on clip organization, role definitions, and coverage of vocal feedback. We further organize Seamless into speaker-held-turn clips with aligned listener audiovisual media, and use timestamped visual and vocal cues to obtain ordered reaction descriptions.

% Paragraph role P06: 既有评价与本文协议
Feedback evaluation has also moved beyond visual quality. REACT2023~\cite{react2023} and ReactFace~\cite{reactface} consider multiple appropriate reactions. GLARE evaluates category matches with temporal overlap. OmniResponse and SpeakerVid-5M assess multimodal response quality and audiovisual synchronization, while OmniPersona also evaluates the appropriateness of listening expressions. Building on these dimensions, our evaluation design separates clip-level event-category agreement, contextual appropriateness and engagement, and coordination within generated audiovisual feedback. Clip-level agreement does not require alignment with one reference's event times, but cannot by itself establish contextual appropriateness. The three dimensions address short, intermittent feedback that may have multiple reasonable realizations.
\cnnote{引用与审阅}{本节为中文相关工作草稿的英文整理，仍待作者审阅。GLARE和OmniPersona沿用提供稿件，不推断作者或录用信息；正式BibTeX按引用清单补齐。}
